% Grid Data: map of the grid representation per country (scripts/plot_grid_map.py), then the standard procedure
% (scripts/plot_grid_steps.py) and the adjustments (scripts/plot_hvdc_check.py, scripts/plot_osm_entries.py, scripts/plot_grid_check.py, scripts/plot_isolated.py).
% The subsubsections are bullets on purpose: material to be turned into prose. The adjustments are one paragraph per commit of the fork.
\subsection{Grid Data}\label{sec:grid}

\begin{figure}[htb]
  \centering
  \includegraphics[width=\textwidth]{grid_map.pdf}
  \caption{How the transmission grid of each modelled country enters the model.
  One node per country: the country is a single bus without an internal grid.
  OpenStreetMap grid: lines and substations from OpenStreetMap, processed by PyPSA-Earth into a network of 50 buses (Section~\ref{sec:grid-standard}), with the corrections of Section~\ref{sec:grid-adjust}; where marked, small grid fragments without a connection to the main grid are also connected to its nearest bus.
  Country borders from \citet{naturalearth}.}
  \label{fig:gridmap}
\end{figure}

For the archetype country US, Brazil, North-West Europe, China and India the spatial ditribution of the power system, and therefore the grid, is modelled (\ref{fig:gridmap}).
The 116 countries with one node skip this processing step.
Of the spatially resolved regions, the (contiguous) United States and Brazil take their grid from the PyPSA-Earth method, apart from the commits of Section~\ref{sec:grid-adjust}.
China, India and the twelve countries of North-West Europe also change the handling of isolated grid fragments, described in the same section.

\subsubsection{Standard procedure}\label{sec:grid-standard}

\begin{figure}[!tp]
  \centering
  \includegraphics[width=\textwidth]{grid_steps.pdf}
  \caption{The standard procedure for Brazil.
  (a) Power lines and cables in OpenStreetMap, coloured by their highest tagged voltage.
  (b) Base network: the lines of \SI{200}{kV} and above, joined at buses that stand for substations, after cleaning.
  (c) Network of the optimisation: the base network merged into 50 buses, line width proportional to capacity, the region each bus represents in the background.
  The bus without lines in Bahia is the merged node of Brazil's small grid islands (148\,MW mean load).
  (d, e) The same close-up of São Paulo (boxes in a and b) on OpenStreetMap's own map, for direct comparison: (d) the lines and substations in OpenStreetMap, (e) the base network, with its lines drawn straight from bus to bus.
  Map tiles \copyright{} OpenStreetMap contributors.}
  \label{fig:gridsteps}
\end{figure}

\paragraph{Data Exctraction and Processing} (rules \texttt{download\_osm\_data, clean\_osm\_data, build\_osm\_network})

\begin{itemize}[nosep,leftmargin=*]
  \item Grid data is obtained using earth-osm \citep{earthosm} obtaining four OpenStreetMap features per country: \texttt{power=line}, \texttt{cable}, \texttt{substation} and \texttt{generator}.
  \item Lines that are tagged with frequency 0 are classified as High Voltage Direct Current (HVDC).
\end{itemize}

\begin{itemize}[nosep,leftmargin=*]
  \item Only lines with both ends inside the modelled area are kept, and each line gets the country it lies in.
  \item Everything below \SI{200}{kV} is dropped (PyPSA-Earth default \SI{51}{kV}); this step de facto only leaves the transmission level in all countries where we model the power grid.
  \item The São Paulo close-up is an example of this, 570 lines and 305 substations go in, and most of them belong to the city's grid below \SI{200}{kV} (Figure~\ref{fig:gridsteps}d).
\end{itemize}

\begin{itemize}[nosep,leftmargin=*]
  \item OpenStreetMap does not connect lines to subtations and therefore does not offer the actual grid topology, which instead has to be estimated.
  \item This estimation first aggregates substantion that are within 5km of each other to lower the burden on data precision.
  \item Every line end is attached to the nearest bus of the same voltage and kind (AC or DC).
  \item This topolog inference greatly reduces the complexity of the grid, e.g. São Paulo: 15 buses and 21 lines remain of the 305 substations and 570 lines in the close-up (Figure~\ref{fig:gridsteps}d, e).
\end{itemize}

\paragraph{Electrical parameters} (rule \texttt{base\_network}; Figure~\ref{fig:gridsteps}b)
\begin{itemize}[nosep,leftmargin=*]
    \item OpenStreetMap does not contain wire types per line, which are therefore assumed to be a standard conductor type depending on line voltage.
  \item A conductor type is a wire design from PyPSA's line-type table, e.g.\ Al/St 240/40 4-bundle: four aluminium conductors with steel core of \SI{240}{mm^2} per phase, with a thermal current limit and resistance, reactance and capacitance per km.
  \item The type sets the line's thermal capacity; its impedances are replaced by \SI{380}{kV} equivalents in the simplification step.
\end{itemize}

\begin{itemize}[nosep,leftmargin=*]
  \item Its capacity is $s_\text{nom} = \sqrt{3}\, i_\text{nom}\, v\, n_\text{circuits}$, with $i_\text{nom}$ the current limit of one phase of the conductor type (kA), $v$ the line's voltage from OSM (kV, phase to phase), $n_\text{circuits}$ the number of three-phase circuits, and $\sqrt{3}$ from the three phases.
  \item For example, one circuit at \SI{230}{kV} gets $\sqrt{3} \times \SI{1.29}{kA} \times \SI{230}{kV} = \SI{514}{MVA}$, with unit conversion 1 V $\cdot$ A = 1 J/C $\cdot$ 1 C/s = 1 J/s = 1 W and linear approximation is 514MW.
  \item To reflect a security margin, only up to 70\,\% of that capacity can be used ($s_\text{max\_pu} = 0.7$); we limit long circuits further by their loadability and use a margin of 0.4 (Section~\ref{sec:grid-adjust}).
  \item Each HVDC line becomes a PyPSA link component, with its capacity set as described in Section~\ref{sec:grid-adjust}.
  \item Transformers get \SI{2000}{MVA}.
  \item Lines that are tagged \textit{under construction} get zero capacity.
\end{itemize}


\paragraph{Regions} (rule \texttt{build\_bus\_regions})
\begin{itemize}[nosep,leftmargin=*]
  \item Every bus gets the Voronoi cell around it, clipped to the country, onshore and in the exclusive economic zone.
  \item These cells are the basis for renewable potentials, power plants and demand of the other layers.
  \item The resulting stack of per-cell components is the unit that is then aggregated for computations feasibility.
\end{itemize}

\paragraph{Simplification} (rule \texttt{simplify\_network})
\begin{itemize}[nosep,leftmargin=*]
  \item To simplify spatial aggregation, all lines are mapped onto the voltage level \SI{380}{kV} with impedance rescaled to preserve thermal capacity.
  \item Transformers are removed, and each bus of a lower voltage is merged into the higher-voltage bus of its station.
  \item A chain of HVDC links between two branching points becomes one link (Section~\ref{sec:grid-adjust}).
  \item Dead ends (buses with a single connection) are folded into their neighbour.
  \item Isolated AC sub-networks below \SI{20}{MW} mean load are dropped, including their plants.
  \item The remaining sub-networks below \SI{300}{MW} are merged, per country (Figure~\ref{fig:gridsteps}c).
\end{itemize}

\paragraph{Clustering} (rule \texttt{cluster\_network})
\begin{itemize}[nosep,leftmargin=*]
  \item In this step, the OSM-drive numbers of regional nodes are aggregated to a user-defined number which usually complies with the available computational resources, for the purpose of this exposition, we assume this number is 50.
  \item The 50 nodes of a region are shared out over its countries and AC sub-networks in proportion to their mean load.
  \item The step is executed as as a small integer programme and every sub-network gets at least one node.
  \item Within each share, k-means on the bus coordinates places the nodes, with buses weighted by their load and conventional plant capacity.
  \item Lines between two clusters are combined into one: capacities add up, and the length is the straight-line distance between the cluster centres times 1.25.
  \item The combined line so carries the summed thermal capacity of its circuits; its own length, 227 to \SI{520}{km} at the median of a region, is that of none of them (Figure~\ref{fig:linelengths}).
  \item HVDC links are not combined; each keeps its own capacity between its two clusters.
  \item The clustered regions are the unions of the Voronoi cells of their buses.
  \item For Brazil this results in 50 buses, 88 lines and 9 HVDC links Figure~\ref{fig:gridsteps}c.
\end{itemize}

\FloatBarrier
\subsubsection{Adjustments}\label{sec:grid-adjust}

Each adjustment is one commit on the branch \texttt{catalyst} of our fork of PyPSA-Earth, named after what it does and linked to the commit, in the order of the workflow.
Options are off by default in the fork and switched on in the configurations of the five regions, \texttt{dc\_ratings} and \texttt{loadability} in all stages; commits without an option always apply.
Section~\ref{sec:grid-check} checks the result as a whole.

\paragraph{Keep a way once across country extracts} (\commit{4d1355aa})

A cable that crosses a border arrives once per country extract, for example way 117931185 of IFA 2000 in the French and in the British file.
The cleaning step turned the copies into parallel lines by adding a suffix to their identifiers.

Only one row per OpenStreetMap identifier is kept, the same idea as in PyPSA-Eur \citep{horsch2018pypsa}.
The feature becomes active in rule \texttt{clean\_osm\_data}, depending on the option:
\yamlblock{drop_duplicate_ids}

\paragraph{Count circuits, not cables, on ways with several voltages} (\commit{ea56f416})

When a way has several voltages and one total of cables, \texttt{fill\_circuits} writes the number of cables where the number of circuits belongs.
A way tagged \texttt{voltage=230000;230000}, \texttt{cables=6}, \texttt{circuits=2} becomes two lines of three circuits each, six circuits instead of two, and with them three times the capacity.
This concerns 1,404 ways in the US and 670 in North-West Europe: in the 1,061 and 569 of them with a \texttt{circuits} tag, the upstream code exceeds the tag in 99\,\% and 96\,\% (Table~\ref{tab:circuits}).

Such a way now gets its circuits per voltage, in the example one per line and two in total.
The model then exceeds the tag in 5 and 9 ways.
It is the largest error found: AC circuit km fall by 21\,\% in the US and 22\,\% in North-West Europe with this commit and the one above, and by under 1\,\% in Brazil, India and China.
The feature becomes active in rule \texttt{clean\_osm\_data}, depending on the option:
\yamlblock{fix_indirect_circuits}

\begin{table}[htb]
  \centering\small
  \caption{Real ways (a \textit{Way} is an OSM data type consisting of an ordered list of nodes forming a polyline or closed shape) with several voltages and one total of cables, and the circuits per voltage they get without and with the commit.
  The first is the Grand Coulee--Schultz line, tagged with two names and two circuits.
  Lines below \SI{200}{kV} are dropped afterwards, so the values at \SI{115}{kV} and \SI{110}{kV} do not reach the network.
  NWE: North-West Europe.
  ``Ways alike'': ways of the base network in that region with the same number of voltages, cables and \texttt{circuits} tag.}
  \label{tab:circuits}
  \input{../build/circuit_ways.tex}
\end{table}

\paragraph{Measure line lengths on the ellipsoid} (\commit{1b5be765})

\textit{What needed fixing.}
Lengths were measured in Web Mercator (EPSG:3857), which stretches them by $1/\cos(\text{latitude})$.
On HVDC lines that is a factor of 1.05 in Brazil, 1.08 in India, 1.18 in China, 1.33 in the US and 1.67 in North-West Europe.
HVDC links keep this length through clustering, so it enters their costs; AC line lengths are recomputed in the clustering.

\textit{What was done.}
Lengths are now measured on the WGS84 ellipsoid.
The feature becomes active in rule \texttt{build\_osm\_network}, depending on the option:
\yamlblock{geodesic_length}

\paragraph{Give HVDC lines a rated capacity} (\commit{1e98e55e})

\textit{What needed fixing.}
PyPSA-Earth reads only the voltage of an HVDC line from OpenStreetMap and gives every line the current of one \SI{500}{kV} cable (\SI{1.638}{kA}), so $p_\text{nom} = v \times \SI{1.638}{kA}$ and \SI{1.31}{GW} at $\pm$\SI{800}{kV}, where Chinese and Indian bipoles carry 6--\SI{8}{GW}.

\textit{What was done.}
A table holds ratings per country and pole voltage, the median design rating of all projects in Wikipedia's list of HVDC projects that are not decommissioned, plus Changji--Guquan ($\pm$\SI{1100}{kV}, \SI{12}{GW}), which the list lacks.
A line takes the row of its country within 10\,\% of its voltage, else the world median of the nearest voltage.
Converters take the summed rating of the HVDC lines they serve, instead of a flat \SI{2000}{MW}.
In the five regions the next commit replaces these ratings.
The feature becomes active in rule \texttt{base\_network}, depending on the option:
\yamlblock{dc_ratings}

\paragraph{Assign HVDC capacities per project, matched by way name} (\commit{38420f22})

\textit{What needed fixing.}
One rating per country and voltage keeps the rule ``one OpenStreetMap line is one bipole'' and guesses the rating instead of looking the project up.
The France--Great Britain border is the clearest case, because its real capacity is not in doubt (Table~\ref{tab:hvdcfrgb}, Figure~\ref{fig:hvdccheck}a).

\begin{table}[htb]
  \centering\small
  \caption{France--Great Britain HVDC: the three operating projects against one rating per country and voltage, in the clustered network.}
  \label{tab:hvdcfrgb}
  \begin{tabular}{@{}lrrp{0.5\textwidth}@{}}
    \toprule
    Project & Real, GW & Median rule, GW & What happens \\
    \midrule
    IFA 2000 ($\pm$\SI{270}{kV}) & 2.0 & 4.0 & drawn as parallel ways, each given the full \SI{2}{GW} of the project; the two landward ways set the limit \\
    ElecLink ($\pm$\SI{320}{kV}) & 1.0 & 2.4 & the link appears twice, each with the \SI{1.2}{GW} median of Great Britain at that voltage \\
    IFA-2 ($\pm$\SI{320}{kV}) & 1.0 & 1.0 & correct \\
    \midrule
    Total & 4.0 & 7.4 & \\
    \bottomrule
  \end{tabular}
\end{table}

\begin{itemize}[nosep,leftmargin=*]
  \item Separate ways: OpenStreetMap may draw each pole, cable pair or bipole as its own way, and each way then gets the rating of a whole project; German SuedLink, for example, has ways named \emph{Plus} and \emph{Minus}.
  \item Lines not yet in service: the OpenStreetMap snapshot is from September 2026 and the fleet is calibrated to 2024; the Bornholm Energy Island cable, under construction, accounts for \SI{4}{GW} of the \SI{4.7}{GW} between Germany and Denmark.
  \item Lines missing: BritNed and COBRA are in OpenStreetMap but almost none of their ways reach the base network, for a reason not yet traced, and Kontek and the Great Belt link are not found under their names, so these borders carry no capacity.
  \item Medians: a correct line can still get another project's rating, such as ElecLink's \SI{1.2}{GW} instead of \SI{1.0}{GW}.
  \item Stale list: the HVDC list shows many Chinese $\pm$\SI{800}{kV} lines at \SI{6.4}{GW} and under planning names, whereas the 20 UHVDC lines in service by 2024 total \SI{158.6}{GW}, with a typical rating of \SI{8}{GW} \citep{wikiuhv}.
\end{itemize}

\textit{What was done.}
OpenStreetMap says where a line runs, and a project table how much it carries.
The feature becomes active in rule \texttt{base\_network}, depending on the option:
\yamlblock{dc_projects}
\begin{itemize}[nosep,leftmargin=*]
  \item The ways hold more than the voltage: names on 91--100\,\% of HVDC km in every region, and ratings on 63\,\% of it in North-West Europe but almost none elsewhere; earth-osm keeps most of these in a column of rare tags (\texttt{other\_tags}), which this commit reads.
  \item The table takes ratings, lengths and years from the Wikipedia list, coordinates also from PyPSA-Eur's copy of it \citep{horsch2018pypsa}, and hand-curated corrections with a source per row, including all Chinese UHVDC lines.
  \item A second hand-curated table links every OpenStreetMap name to its project, because the two often differ: \emph{HVDC Lingzhou - Shaoxing} is the list's Yinchuan--Zhuji, and some names are only in Chinese.
  \item Matching by name comes first (Figure~\ref{fig:osmentries}a): a project takes the ways carrying its names, and its terminals pick one route through parallel ways, so that it follows its own line.
  \item If the named ways do not connect its terminals, the route continues through unnamed ways (Figure~\ref{fig:osmentries}b), preferring the named ones.
  \item Without terminals or a route, the project takes all its named ways, and if those cover under half its listed length, or it has none, it counts as missing.
  \item Every bus pair on a route carries the project rating once, shared by its parallel links, so the number of ways no longer matters.
  \item Projects entering service after 2024 get zero capacity, as do links matching no project.
  \item A project in service with no usable route becomes one link between the AC buses of its countries at its terminals: BritNed, COBRA, Kontek and the Great Belt link in North-West Europe, and Hubei--Shanghai, Baihetan--Jiangsu, Baihetan--Zhejiang and two small links in China.
  \item Converters again take the summed rating of their lines.
\end{itemize}

\begin{figure}[H]
  \centering
  \includegraphics[width=\textwidth]{hvdc_check.pdf}
  \caption{(a) The ways of the three France--Great Britain projects in OpenStreetMap.
  (b) HVDC capacity per pair of countries in North-West Europe: the model's network before and after the corrections of this section, and the projects in service by 2024 in the list of HVDC projects \citep{wikihvdc}.
  Greenlink (Great Britain--Ireland, \SI{0.5}{GW}) entered service in 2025 and is excluded.
  The remainder between Germany and Denmark is exactly SylWin1 (\SI{0.864}{GW}), a German offshore wind link whose offshore bus falls in the Danish zone.}
  \label{fig:hvdccheck}
\end{figure}

\begin{figure}[H]
  \centering
  \includegraphics[width=\textwidth]{osm_entries.pdf}
  \caption{Two $\pm$\SI{800}{kV} ways in China's OpenStreetMap data, with their tags, each highlighted among the country's HVDC lines.
  (a) A way whose name can be looked up in the project table; panel (a) also highlights the second way of that line.
  (b) A way without a name, tagged only as an unidentified line; it can be placed by voltage and geometry alone.
  Tags in Chinese script are not shown.}
  \label{fig:osmentries}
\end{figure}

\paragraph{Keep the capacity and length of a merged HVDC chain} (\commit{11f2f5f7})

\textit{What needed fixing.}
When the simplification merged a chain of HVDC segments, the upstream code summed the capacity and length of all segments and converters, so a link's capacity grew with the number of OpenStreetMap segments it was drawn in (one Chinese link reached \SI{29.8}{GW}).

\textit{What was done.}
Parallel segments add up, the weakest group limits the chain, and lengths add along the series.
Rule \texttt{simplify\_network}.

\paragraph{Absorb DC buses left after link simplification} (\commit{caa0f116})

\textit{What needed fixing.}
OpenStreetMap leaves DC buses that the merging of chains cannot fold: converter stations that also carry load or generators, cable stubs ending at a DC bus without a converter, and isolated substations tagged DC.
Each is its own sub-network and claims a node in the clustering, 18 of 50 in North-West Europe.

\textit{What was done.}
Every remaining DC bus is mapped onto the AC bus of its converter, or else the nearest AC bus of its country.
Converters that become self-loops are dropped, and the DC links then run between AC buses.
Rule \texttt{simplify\_network}.

\paragraph{Do not merge buses of different countries through transformers} (\commit{3d9ed567})

\textit{What needed fixing.}
Mapping all lines onto one voltage level merges the two buses of each transformer.
At a border substation these can lie in two countries (Aldeadávila: a Portuguese \SI{220}{kV} and a Spanish \SI{400}{kV} bus \SI{100}{m} apart), so the low-voltage bus and everything on it changed country, which in a one-node run moved Portugal's \SI{56}{TWh} of load into Spain.

\textit{What was done.}
Transformers between countries are kept as \SI{1}{km} lines of the transformer's rating, and the buses stay in their countries.
Rule \texttt{simplify\_network}.

\paragraph{Never treat the backbone as an island} (\commit{d77128a8})

\textit{What needed fixing.}
The join of isolated fragments, an upstream option (\texttt{s\_\allowbreak threshold\_\allowbreak fetch\_\allowbreak isolated}), selects every AC bus whose sub-network holds less than the threshold of national load and removes its lines.
It did not exclude the country's largest sub-network, so a threshold above the backbone's own share left every bus isolated: Russia at a threshold of 1.0 came out as 1,946 one-bus sub-networks.

\textit{What was done.}
The backbone is excluded from the selection.
Rule \texttt{simplify\_network}.

\paragraph{Tolerate an empty lines table} (\commit{064d1990})

\textit{What needed fixing.}
A country with no OpenStreetMap line above the voltage threshold imports zero lines, and the workflow then failed on a missing column; Singapore's cables are underground and unmapped.

\textit{What was done.}
Both accesses to the line table are guarded.
Rule \texttt{base\_network}.

\paragraph{Join isolated fragments to the backbone} (no commit: settings of upstream options)

\begin{figure}[htb]
  \centering
  \includegraphics[width=\textwidth]{isolated.pdf}
  \caption{Connected AC networks (sub-networks) in the base networks of India, China and North-West Europe, and what the simplification step does with them.
  Grey: the main grid of each country, its largest connected network (the backbone).
  Coloured: the smaller fragments, by what happens to them; a fragment is kept as a separate network if it carries 5\,\% of the national load or more, or crosses a border.}
  \label{fig:isolated}
\end{figure}

\textit{What needed fixing.}
\begin{itemize}[nosep,leftmargin=*]
  \item The grids in OpenStreetMap break into many AC sub-networks: 149 in India, 175 in China and 29 in North-West Europe (Figure~\ref{fig:isolated}).
  \item Most are \SI{220}{kV} pockets whose connection runs below the \SI{200}{kV} threshold or is not mapped.
  \item An area tied to the rest only by HVDC also counts as its own AC sub-network.
  \item Under the default rules every fragment above \SI{300}{MW} stays an isolated sub-network and claims one of the 50 nodes, as does the merged node of the smaller ones.
  \item Because these thresholds are absolute, the topology changes with the demand of the scenario: with 2050 demand, India had 54 isolated sub-networks above the thresholds, more than its 50 nodes, and the clustering failed.
  \item The default merge also runs before any join: India's 118 fragments below \SI{300}{MW} (142 buses) would become one node of \SI{12.0}{GW} at their mean position in central India (\ang{78.9}\,E \ang{20.5}\,N), which at 5.2\,\% of India's load stays isolated.
  \item In China the merged node (\SI{13.7}{GW}, 1.2\,\%) would be joined, but onto a single backbone bus in central China, far from where most of that demand is.
\end{itemize}

\textit{What was done.}
\begin{itemize}[nosep,leftmargin=*]
  \item China, India and North-West Europe set:
  \yamlblock{isolated}
  \item The merge is off: fragments below \SI{20}{MW} are still dropped, and every other sub-network of a country except its backbone, below 5\,\% of the country's load, is joined to the backbone, each of its buses to the nearest backbone bus.
  \item No new line is built for this: the fragments are treated as part of the bus they are joined to.
  \item The share is relative to the country's load, so the scenarios with today's and with 2050 demand get the same topology.
  \item Result: 131 of 149 sub-networks are joined and 17 dropped in India, 159 of 175 joined and 15 dropped in China, and 24 of 29 joined and 2 dropped in North-West Europe.
  \item The largest joined fragment holds 0.45\,\% of national load in India and 0.22\,\% in China; in North-West Europe the largest are in Denmark (3.4\,\%) and Ireland (2.9\,\%).
  \item Northern Ireland stays a sub-network of its own, because it is tied by AC to Ireland, and a sub-network spanning two countries is never joined.
  \item India's 118 small fragments now spread over 124 backbone buses, the largest receiving \SI{0.4}{GW}, and China's 122 over 163 buses, the largest receiving \SI{0.3}{GW}.
\end{itemize}

\paragraph{Limit long circuits by their loadability} (\commit{e3b61211})

\begin{figure}[htb]
  \centering
  \includegraphics[width=\textwidth]{line_lengths.pdf}
  \caption{AC transmission capacity between clusters by the length of the circuit, in the six spatially resolved regions, sorted by the median.
  Each row: the circuits of the base network whose two ends lie in different clusters, with their length from OpenStreetMap, as share of their summed thermal rating per length bin.
  Black line: median by capacity.
  Grey area: lengths at which the loadability of a \SI{380}{kV} circuit lies below its thermal rating.
  Left: share of the capacity that the loadability limits.}
  \label{fig:linelengths}
\end{figure}

\textit{What needed fixing.}
A line's capacity is the thermal rating of its conductors, whatever its length.
Heat limits short lines only: a longer line is limited by the voltage drop along it and, longer still, by the stability of the angle between its two ends \citep{stclair1953,dunlop1979}.
The St.~Clair curve expresses this as the load a line without compensation can carry in multiples of its surge impedance loading (SIL): 3.0 at 50 miles and 1.0 at 300 miles (Figure~\ref{fig:loadcurve}a).
The model's line types have a thermal rating of 2.5 to 2.8 times their SIL, so the curve lies below the thermal rating from \SI{93}{km} on at \SI{380}{kV} and from \SI{118}{km} on at \SI{750}{kV}.
The regions differ in how much of their grid is that long (Figure~\ref{fig:linelengths}).
Half of the capacity between clusters lies in circuits longer than \SI{220}{km} in Brazil and \SI{158}{km} in India, against \SI{120}{km} in US West, \SI{64}{km} in US East and China and \SI{38}{km} in North-West Europe.
Brazil and India are also the regions with the most generous grids: at 0.7 of the thermal rating their line limits sum to 4.0 and 5.2 times the peak load, against 0.6 to 1.7 times in the other regions (Figure~\ref{fig:loadcheck}a).

\begin{figure}[htb]
  \centering
  \includegraphics[width=\textwidth]{loadability_curve.pdf}
  \caption{The loadability limit.
  (a) The St.~Clair curve of a line without compensation, in multiples of the surge impedance loading; dots are the points the model interpolates between (3.0 at 50 miles and 1.0 at 300 miles as published \citep{stclair1953}, the others read off the curve as reproduced by \citet{sekhon1981}).
  Horizontal: the thermal rating of two of the model's line types, which is the limit left of the crossing.
  (b) Share of the thermal rating that is loadable, per clustered AC line, against the mean length of the circuits it combines, dot area by thermal rating; grey: one \SI{380}{kV} circuit of that length.}
  \label{fig:loadcurve}
\end{figure}

\textit{What was done.}
Every AC circuit of the base network is limited to $\min(s_\text{nom},\, k(L)\,\mathrm{SIL})$, with $L$ its own length.
$k(L)$ is the St.~Clair curve, interpolated between its points on logarithmic axes and continued with its first slope below 50 miles and with its last slope beyond 600 miles.
Per circuit, $\mathrm{SIL} = v^2 / Z_c$ with the surge impedance $Z_c = \sqrt{x / (\omega c)}$ from the reactance $x$ and capacitance $c$ per km of the line type: \SI{606}{MW} for a \SI{380}{kV} circuit against a thermal rating of \SI{1698}{MVA}, and \SI{2175}{MW} against \SI{5404}{MVA} at \SI{750}{kV}.
The base network is used because there every circuit still has its own length, voltage and line type.
The usable share $s_\text{max\_pu}$ of a clustered line is multiplied by the sum of these limits over the circuits between its two clusters, divided by the sum of their thermal ratings (Figure~\ref{fig:loadcurve}b).
The length of the clustered line, the distance between two cluster centres, does not enter, so the result does not depend on the number of clusters.
The security margin so applies to whichever of the two limits binds.
Capacity that the optimisation adds to a line is limited in the same way, so a corridor of long circuits costs more per usable megawatt than one of short circuits.
HVDC links are not touched.
The loadability limits 84\,\% of the capacity between clusters in Brazil, 68\,\% in India, 45\,\% in US West, 25\,\% in US East, 14\,\% in China and 11\,\% in North-West Europe.
It lowers the line limits by 23\,\% in Brazil, 20\,\% in India, 13\,\% in US West, 5\,\% in US East and 1 to 3\,\% in China and North-West Europe.
The curve is that of a line without compensation \citep{sekhon1981}; series compensation raises what real long lines carry, and the model does not represent it.
The feature becomes active in rule \texttt{prepare\_network}, depending on the option \texttt{lines: loadability}, shown below with the margin.

\textit{A lower security margin.}
The loadability alone does not bring the regions down to their published transfer capacities (Section~\ref{sec:grid-check}): across the borders with a published figure the circuits still carry 1.2 to 2.6 times that figure.
PyPSA-Earth's 0.7 is a margin per line.
Published transfer limits are limits of the whole system: after the loss of any element, with loop flows, with bottlenecks behind the border and under stability criteria.
For each region we computed the margin at which the circuits across its borders, limited by their loadability, carry the published capacity: 0.60 in US West, 0.41 in US East, 0.40 in Brazil, 0.35 in India, and in North-West Europe 0.26 against the net transfer capacities of 2024 and 0.39 against those of the reference grid of 2030.
We use 0.4 for every region, China included, for which no published figure with a border was found.
\yamlblock{loadability}

\begin{figure}[htb]
  \centering
  \includegraphics[width=\textwidth]{loadability_map.pdf}
  \caption{The clustered AC networks with the loadability limit and the security margin of 0.4.
  Line width: thermal rating.
  Colour: usable share of it, 0.7 for every line in PyPSA-Earth.
  The share in each title is weighted by capacity.
  HVDC links are not drawn.}
  \label{fig:loadmap}
\end{figure}

\paragraph{Limit the build-out of the grid} (\commit{f374de79})

\begin{figure}[htb]
  \centering
  \includegraphics[width=\textwidth]{grid_expansion.pdf}
  \caption{The transmission grid relative to 2024 as the model may build it.
  Solid: the cap with the rate at which the region's grid has grown, the central case.
  Dashed: the cap with the rate that plans and needs studies ask for.
  Grey: the growth of the last decade.
  Circles: published plan figures, in circuit km or interface capacity (Table~\ref{tab:gridexpansion} names them).}
  \label{fig:gridexpansion}
\end{figure}

\textit{What needed fixing.}
The scenarios let the optimisation reinforce every AC corridor and every HVDC link at cost (\texttt{ll: copt}), without a bound, a lead time or a build rate.
Real grids grow slowly: new transmission takes five to fifteen years to plan, permit and build \citep{ieagrids2023}.
With lines limited by their loadability and the lower security margin, India's zero-carbon run answered by adding \SI{309}{GW} of AC line capacity and \SI{22}{GW} of HVDC in one step, and China's by adding \SI{1900}{GW} of HVDC to the \SI{354}{GW} it has.
An unbounded grid makes distant resources too cheap and local firm capacity too small.

\textit{What was done.}
The transmission volume, capacity times length summed over AC lines and HVDC links, is limited to that of the existing grid grown at a rate per year: $V(y) \le (1 + r)^{\,y - 2024}\, V_{2024}$.
The limit is one global constraint per region, so the optimisation chooses where the grid grows, and HVDC counts in the same limit as AC.
In the pathway it is set anew for every planning horizon, and what was built in a horizon stays as the minimum of the next.
The single-step runs of 2050 get the limit at its 2050 value; the dispatch of 2024 needs none, because its grid is frozen (\texttt{ll: v1.0}).
The rate $r$ of a region is the growth its grid has delivered (Table~\ref{tab:gridexpansion}, Figure~\ref{fig:gridexpansion}): 1\,\% per year in the two US regions and North-West Europe, 2\,\% in Brazil, 3\,\% in India and 4\,\% in China.
The rates that official plans and needs studies imply are kept as a sensitivity.
Published growth is in circuit km or in interface capacity, not in capacity times length; we assume the volume grows like them.
Uprating of existing lines, common in the US and Europe, is not in the length figures.

\begin{table}[htb]
  \centering
  \caption{The cap on grid build-out: rate per year and the grid it allows relative to 2024.
  Delivered: US about 1\,\% per year in the last decade and 2\,\% from 1978 to 2020 \citep{repeat2022}; EU $+12$\,\% in a decade, Brazil more than $+50$\,\%, India about $+60$\,\% \citep{ieagrids2023}; India $+3$\,\% per year from 2019-20 to 2024-25 \citep{powerline2025}; China's west-to-east capacity from 270 to \SI{340}{GW} between 2020 and 2025.
  Plans: US 2.1 to 2.6 times the system of 2020 by 2050 \citep{gridstrategies2025}; cross-border capacity between the twelve European countries from 46.6 to \SI{71.3}{GW} by 2030 \citep{emberinterconnection}; Brazil 191,353 to \SI{220134}{km} by 2035 \citep{epepde2035}; India 496,372 to 648,190 circuit km by 2031-32 \citep{powerline2025}; China \SI{420}{GW} by 2030.}
  \label{tab:gridexpansion}
  \small
  \input{../build/grid_expansion.tex}
\end{table}

The feature becomes active in rules \texttt{add\_brownfield} and \texttt{prepare\_network} with the option \texttt{lines: expansion\_limit}.
The driver writes the option for every stage from the table of rates, by the column named in the first block; the second block gives the single-step runs of 2050 their year:
\yamlblock{expansion_limit}


\FloatBarrier
\subsubsection{Check against real project metrics}\label{sec:grid-check}

The script \texttt{check\_grid.py} applies the corrections in memory to the base networks and compares them with the project table.
Summing link capacities is not a valid test, because a project crossing several clusters counts once per link, which is why China's \SI{467}{GW} against about \SI{175}{GW} on the Wikipedia list says little about transfer capacity.
Capacity times true length avoids this (Table~\ref{tab:hvdcregions}, Figure~\ref{fig:gridcheck}a).

\begin{table}[htb]
  \centering\small
  \caption{HVDC in the base networks against the listed projects in service by 2024, with their listed lengths.
  ``Uncorrected'': one rating per country and voltage, with Web Mercator lengths.
  ``Projects'': in service and found in the network, plus those added as new links, of all listed, including small links below the voltage threshold.
  ``Routes'': matched projects whose route length lies within 10\,\% of the listed length.}
  \label{tab:hvdcregions}
  \begin{tabular}{@{}lrrrrrr@{}}
    \toprule
    & \multicolumn{3}{c}{GW\,km, thousand} & & & \\
    \cmidrule(lr){2-4}
    Region & Listed & Uncorrected & Corrected & Corrected / listed & Projects & Routes \\
    \midrule
    US & 8.0 & 12.1 & 7.9 & 0.98 & 7 of 7 & 6 of 7 \\
    Brazil & 38.5 & 40.4 & 38.7 & 1.01 & 6 of 6 & 6 of 6 \\
    India & 39.9 & 44.6 & 39.4 & 0.99 & 9 of 9 & 9 of 9 \\
    China & 330.2 & 356.0 & 314.4 & 0.95 & 32 + 5 of 40 & 25 of 32 \\
    North-West Europe & 5.4 & 19.7 & 4.2 & 0.76 & 20 + 4 of 27 & 7 of 20 \\
    \bottomrule
  \end{tabular}
\end{table}

\begin{figure}[H]
  \centering
  \includegraphics[width=\textwidth]{grid_check.pdf}
  \caption{The model's grid before and after the corrections of this section, against real metrics.
  (a) HVDC capacity times length relative to that of the operating projects in the list of HVDC projects (vertical line).
  (b) AC circuit km against published totals: India May 2024, 220 kV and above without HVDC \citep{indiatransmission}, Germany about 37,000 km \citep{bmwkgrid}, Brazil 2020 with the diamond at the 28\,\% growth projected for 2024 \citep{ipeabrazil}.
  ``Before correction'' carries the line lengths measured in the Web Mercator projection, which overstate distances away from the equator; the lines of the optimisation do not use them.}
  \label{fig:gridcheck}
\end{figure}

\begin{itemize}[nosep,leftmargin=*]
  \item Matching by name replaces an earlier version that used terminals alone: that one reached the right national totals, but put about a quarter of China's named km on the wrong parallel line, for example Baihetan--Jiangsu on ways named \emph{HVDC Jinping--Sunan}, with a length close enough that the route test did not notice.
  \item By construction the name now agrees, so the remaining tests are route length, capacity times length, borders and rating tags, none of which the name match uses.
  \item Route lengths hold where lines are mapped whole: all projects in Brazil and India, and for China \SI{3302}{km} against \SI{3293}{km} for Changji--Guquan, with 24 of 27 fully named lines within 10\,\%.
  \item The US miss is the partial Quebec--New England section, whose length we took from OpenStreetMap, and we could not confirm that the way named \emph{Nexus} is the CU project, although its length agrees.
  \item China's total lies 5\,\% below the list, mainly because four lines are only partly named or connected (65--86\,\% of their listed length), so their other sections are not counted: Northwest Yunnan--Guangdong, Wudongde, Yazhong--Jiangxi and Shanbei--Hubei.
  \item Seven of eight cross-border pairs in North-West Europe come within 10\,\% of the list (Figure~\ref{fig:hvdccheck}b), and France--Great Britain falls from \SI{7.4}{GW} to \SI{4.0}{GW}.
  \item North-West Europe's total stays low, and its route test fails, mainly because of German and British offshore wind links, whose export cables are only partly in OpenStreetMap, with two more (Dogger Bank A and B) dropped for that reason.
  \item Where a way carries a rating, a tag the matching does not use, it agrees with the table on 22 of 26 links; the others are single cables of Moyle (\SI{225}{MW}, table \SI{500}{MW}) and the Spittal section of the Caithness link (\SI{800}{MW}, table \SI{1200}{MW}).
  \item Unmatched links get zero capacity: 107 of 518 in China and 48 of 133 in North-West Europe, largely unnamed pieces of lines added separately, later lines and links outside the table; a real line missing from the table would be lost.
  \item Series links remain: a matched project still touches a median of 3 clusters in China (at most 7), so each cluster on its way can exchange power with a line that has no converter station there; this is not fixed, and it does not change the capacity along the line.
\end{itemize}

\textbf{AC lines.}
\begin{itemize}[nosep,leftmargin=*]
  \item The HVDC defect does not carry over, because an AC line's capacity is computed per circuit and not taken from a project; the circuit defect above is its counterpart.
  \item Published totals give a first, partial test (Figure~\ref{fig:gridcheck}b): India 461,800 against 466,560 circuit km, with 206,100 against 203,890 at \SI{400}{kV} and 59,100 against 54,797 at \SI{765}{kV}.
  \item Germany comes to 32,600 km against about 37,000 km, against 79,800 km uncorrected, 50,800 km of them with true lengths.
  \item Brazil's 173,500 km compares with about 157,000 km projected for 2024, from an older source.
  \item The US falls from 476,500 to 295,200 circuit km; we found no national total to compare it with, and none for China.
  \item These totals test line lengths and circuit counts only: thermal ratings still come from standard conductor types.
\end{itemize}

\textbf{Transfer capacity across borders.}
\begin{figure}[H]
  \centering
  \includegraphics[width=\textwidth]{loadability_check.pdf}
  \caption{The line limits against real metrics.
  (a) Sum of the limits of the clustered AC lines per GW of peak load, at PyPSA-Earth's 0.7 of the thermal rating and with the loadability limit and the security margin of 0.4.
  (b) Capacity across the borders with a published transfer capacity, model over published: the AC circuits of the base network whose two ends lie on either side of the border, with the HVDC links across it, summed over the borders of a region.
  Published: the United States \citep{nercitcs2024}, Brazil \citep{epepde2035}, India \citep{powerline2025}, Europe \citep{emberinterconnection}, for 2024 and for the reference grid of 2030.
  China has no published figure with a border.}
  \label{fig:loadcheck}
\end{figure}

\begin{itemize}[nosep,leftmargin=*]
  \item The script \texttt{check\_loadability.py} calls the rule's function in memory on the networks, so Figures~\ref{fig:linelengths} to~\ref{fig:loadcheck} do not depend on a solved network.
  \item It counts the AC circuits of the base network whose two ends lie on either side of a border, by the state or country of each end, so the count does not depend on the clustering.
  \item United States: NERC's transfer capability study gives the total transfer capability between neighbouring planning regions for the summer and winter peak of 2024 \citep{nercitcs2024}; nine borders in the West and five in the East follow state borders and are compared, each with the larger of its directions and seasons.
  \item The study computes one border at a time, so its values are not simultaneous, and it draws its regions on a map: which states make up a region is our reading of that map.
  \item Brazil: the mean interchange limits of the existing system of September 2025, \SI{13.5}{GW} of export from the Nordeste and \SI{11.0}{GW} of import into the Sul \citep{epepde2035}.
  \item India: \SI{118.7}{GW} between its five regional grids at the end of 2024-25 \citep{powerline2025}; the source does not say whether HVDC is counted, and we take it to be.
  \item Europe: the net transfer capacities of 2024 per border \citep{emberinterconnection}, 16 borders between the twelve countries with AC circuits.
  \item At 0.7 of the thermal rating the model holds 1.3 times the published capacity in US West, 1.7 times in US East, 2.3 times in Brazil, 2.1 times in India and 2.6 times in North-West Europe (Figure~\ref{fig:loadcheck}b).
  \item The loadability brings that to 1.2 in US West, 1.8 in Brazil and 1.7 in India and leaves US East and Europe as they are, because their border circuits are short (\SI{20}{km} and \SI{17}{km} at the median against 212 and \SI{178}{km} in Brazil and India).
  \item With the margin of 0.4 the model holds 0.7 times the published capacity in US West, 1.0 times in US East and Brazil, 1.1 times in India and 1.5 times in North-West Europe.
  \item US West is compared with the larger of two directions and two seasons per border, so 0.7 is against the upper end of the published values.
  \item North-West Europe's net transfer capacities of 2024 are commercial values that their source calls a conservative picture of the physical capacity; against the reference grid of 2030 on the same borders (\SI{56}{GW} instead of \SI{38}{GW}) the model holds 1.0 times the published capacity.
  \item A cluster that spans a border holds the circuits across it inside, where the clustered network has no line and therefore no limit: \SI{128}{GW} of the \SI{312}{GW} between India's regional grids.
\end{itemize}

\textbf{The cap on build-out.}
\begin{itemize}[nosep,leftmargin=*]
  \item India's zero-carbon run of 2050 may build 2.16 times the existing grid and builds 1.16 times: the cap does not bind, and the result equals that without it.
  \item North-West Europe's may build 1.295 times and builds exactly that, with a shadow price of 436 \$ per MW\,km; inside the allowance the AC volume grows by 14\,\% and the HVDC volume more than sixfold.
  \item In a pathway of a ten-node test region with 2\,\% per year, the allowance is 1.02, 1.13 and 1.24 times the existing grid in 2025, 2030 and 2035, and the grid built is never above it and never below that of the horizon before.
\end{itemize}

\textbf{What stays unvalidated.}
\begin{itemize}[nosep,leftmargin=*]
  \item Any effect on model results: load shedding, flows, prices and the technology mix.
  \item The cap on build-out in a pathway of a full region, and with the rates of the plans.
  \item AC transfer capacity in China: no published figure with a border was found (\SI{340}{GW} of west-to-east transmission, AC and HVDC not separated).
  \item The security margin: one value from 33 borders in five regions, applied to every line of every region; China's has no published figure behind it.
  \item The loadability curve is a planning curve for single lines without compensation; circuit lengths are those between two substations in OpenStreetMap, which can be shorter than the section between two points with voltage support.
  \item The corridors of the Chinese projects that are only partly named or were added as straight links.
  \item The hand-curated name table, including our reading of the Chinese names, and the two projects added from memory (Yongren--Funing, the Quebec--New England section).
  \item The project table remains mostly one source, which also supplies the ratings, so the comparison tests the mapping and not the ratings themselves.
\end{itemize}

\subsubsection{One-node countries}

\begin{itemize}[nosep,leftmargin=*]
  \item The same raw data, cleaning and topology steps, with a threshold of \SI{100}{kV} (Canada and Russia \SI{200}{kV}), so that every country has at least one bus.
  \item The region of every bus is the whole country (GADM level 0), and clustering maps all buses of a country onto one.
  \item Isolated buses are kept with their plants instead of being dropped.
  \item All lines and links between countries are removed before the optimisation, so countries do not trade.
  \item Canada and Russia are too large for one whole-country region in the renewable potential calculation, so they take Voronoi regions at \SI{200}{kV}, every fragment is joined to the backbone, and the network is clustered to one node.
  \item Singapore's substations are in OpenStreetMap but the underground cables between them are not, so every bus is isolated; all of them are merged into one bus.
\end{itemize}
